All cluster calculations were performed with Gaussian~16\cite{g16}
using the PBE exchange--correlation functional (PBEPBE).\cite{perdew:pbe:1996}
A def2-TZVP basis was used for C, O, and Ge,\cite{weigend:def2:2005}
and Pt was described with the corresponding Karlsruhe valence basis and a
60-electron \textcolor{cyan}{scalar relativistic} effective core potential \textcolor{cyan}{(ECP)}.\cite{andrae:epp:1990,weigend:def2:2005}
This leaves 18 explicit Pt electrons.
Geometry optimizations and harmonic
frequency calculations were carried out for the bare clusters, carbonyl
complexes, and free \ce{CO}. All structures in the adsorption
analysis were confirmed as true minima by the absence of imaginary frequencies.
Orbital stability was tested for the selected triplet
\ce{Pt10} solution and for the seven distinct triplet \ce{Pt10CO} minima
obtained after placing \ce{CO} at independent Pt sites.
All these wavefunctions were stable under the
perturbations considered.

The structural families used in the adsorption study are summarized in
Figure~\ref{fig:neutral_data_sets}. For the crossed-parent models, the spin
states were optimized separately before \ce{CO} adsorption was considered.
The \ce{Pt6Ge4}, \ce{Pt5Ge5}, and \ce{Pt4Ge6} structures were generated from
the reported global-minimum \ce{Pt10} framework.\cite{ugartemendia:2022} In the reverse route, the
triplet \ce{Pt10} structure was generated from the reported global-minimum
\ce{Pt5Ge5} framework. The reported global-minimum structure of each
composition was also reoptimized at the same level and used in a separate
comparison set.

\begin{figure}[H]
\centering
\begin{tikzpicture}[
  x=1cm,
  y=1cm,
  header/.style={font=\sffamily\bfseries\footnotesize, align=center},
  crossbox/.style={draw=black!65, line width=0.65pt, fill=black!3,
    text width=5.45cm, align=center, inner sep=4pt,
    font=\sffamily\footnotesize},
  gmbox/.style={draw=blue!70!black, line width=0.65pt, fill=blue!4,
    text width=5.45cm, align=center, inner sep=4pt,
    font=\sffamily\footnotesize},
  commonbox/.style={draw=black!55, line width=0.65pt, fill=black!2,
    text width=6.22cm, align=center, inner sep=4pt,
    font=\sffamily\footnotesize},
  crossarrow/.style={->, >=stealth, line width=0.65pt, black!75},
  gmarrow/.style={->, >=stealth, line width=0.65pt, blue!70!black, dashed}
]

\node[header] at (-3.25,2.25)
  {Crossed-parent models\\Two complementary routes};
\node[header,text=blue!70!black] at (3.25,2.25)
  {Reported global-minimum models\\Compare parent isomer and spin state};

\node[crossbox] (crossedstart) at (-3.25,0.68)
  {\textbf{Start}\\
   \ce{Pt5Ge5 -> Pt10}; \quad
   \ce{Pt10 -> Pt6Ge4}, \ce{Pt5Ge5}, \ce{Pt4Ge6}\\
   (see Figure~\ref{fig:tfg_crossed_models})};
\node[gmbox] (gmstart) at (3.25,0.68)
  {\textbf{Start}\\
   Reported GM structures\cite{ugartemendia:2022} of \ce{Pt10},
   \ce{Pt6Ge4}, \ce{Pt5Ge5}, and \ce{Pt4Ge6}};

\node[commonbox] (adsorption) at (0,-2.0)
  {\textbf{Common \ce{CO}-adsorption protocol}\\
   \ce{CO} is placed at selected Pt sites; each structure is optimized\\
   at the selected parent geometry and spin state};

\node[crossbox] (crosseduse) at (-3.25,-5.35)
  {\textbf{Purpose}\\
   \ce{Pt10 -> Pt_xGe_{10-x}}: representative composition trend\\
   \ce{Pt5Ge5 -> Pt10}: pure-Pt geometry and spin\\
   Tables~\ref{tab:tfg_adsorption_descriptors}
   and~\ref{tab:tfg_bimetallic_adsorption_descriptors};\\
   filled black symbols};
\node[gmbox] (gmuse) at (3.25,-4.65)
  {\textbf{Purpose}\\
   Compare adsorption across parent isomers and spin states\\
   Table~\ref{tab:tfg_adsorption_descriptors} and Supporting\\
   Information; open blue symbols};

\draw[crossarrow] (crossedstart.south) -- (adsorption.north west);
\draw[gmarrow] (gmstart.south) -- (adsorption.north east);
\draw[crossarrow] (adsorption.south west) -- (crosseduse.north);
\draw[gmarrow] (adsorption.south east) -- (gmuse.north);
\end{tikzpicture}

\caption{Organization of the two cluster data sets used for the
\ce{CO}-adsorption analysis. The crossed-parent set contains the
\ce{Pt10}-to-\ce{Pt_{x}Ge_{10-x}} composition sequence and the reverse
\ce{Pt5Ge5}-to-\ce{Pt10} route. Calculations initiated from the reported
global-minimum structures show how the results change with the parent isomer
and spin state. Both sets follow the same adsorption protocol.}
\label{fig:neutral_data_sets}
\end{figure}

For this DFT series, the reported relative energies and \ce{CO}
binding energies include the harmonic zero-point vibrational energy (ZPVE).
Hereafter, ``carbonyl'' denotes the optimized cluster--\ce{CO} complex. We define
\begin{align}
E_0(X) &= E_{\rm elec}(X)+\mathrm{ZPVE}(X),\\
BE(\ce{CO}) &= E_0(\mathrm{carbonyl})
              -E_0(\mathrm{bare\ cluster})-E_0(\ce{CO}).
\label{eq:neutral_dft_be}
\end{align}
Thus, more negative values correspond to stronger adsorption.
Unless stated otherwise, each binding energy is referenced to free singlet
\ce{CO} and to the optimized bare cluster with the same parent framework and
multiplicity as the corresponding carbonyl.

The bonding patterns were analyzed with the adaptive natural density partitioning tool (AdNDP~2.0).\cite{tkachenko:adndp2:2019}
AdNDP searches included the chemically active valence space, with 10
electrons per Pt atom and four per Ge atom. The Pt 5s and 5p semicore orbitals
retained by the ECP were treated as core orbitals.
For the triplet \ce{Pt10}, the $\alpha$ and $\beta$ density matrices were analyzed
independently. Elements involving the same atomic centers in both spin
channels were paired when possible. The Occupation Numbers (ON) indicate the probability of finding two electrons in a given bond. The spin-resolved elements and their
ON are reported in Supporting
Table~\ref{tab:tfg_adndp_summary}.

For selected strongly and weakly bound carbonyl minima of \ce{Pt10},
\ce{Pt6Ge4}, \ce{Pt5Ge5}, and \ce{Pt4Ge6}, interaction and
fragment-reorganization terms were evaluated from electronic single-point
energies. The bare-cluster and \ce{CO} fragments were calculated both in their
isolated optimized geometries and at the coordinates they adopt in the
optimized carbonyl complex. The ZPVE correction was added only after this
electronic decomposition, to obtain the final binding energy.
